% !TeX root = ../main.tex

\chapter{绪论}

\section{研究背景}

近年来，深度学习（Deep Learning, DL）技术~\cite{lecun2015deep}在图像识别~\cite{krizhevsky2012imagenet,dosovitskiy2020image}、自然语言处理~\cite{sutskever2014sequence,bahdanau2014neural}等领域取得了重大突破，
推动了以 Transformer~\cite{vaswani2017attention}为代表的基础模型架构的广泛应用。
这一技术演进的一个显著特征是模型参数规模的指数级增长：
在视觉领域，相较于包含约 $2.5 \times 10^7$ 个参数的经典骨干网络 ResNet-50~\cite{he2016deep}，Vision Transformer (ViT-Huge) 的参数量已跃
升至约 $6.3 \times 10^8$~\cite{dosovitskiy2020image}；在自然语言处理领域，BERT-Base 的参数量约为 $1.1 \times 10^8$~\cite{devlin2019bert}，
而 GPT-3 175B 的参数量则达到 $1.75 \times 10^{11}$，在数年内实现了逾三个数量级的跨越式增长~\cite{brown2020language}。模型参数规模
的急剧增长与训练数据量的持续扩大，对底层算力基础设施带来了前所未有的挑战。为此，如图~\ref{fig:datacenter}所示，
企业与研究机构普遍构建了由海量异构加速硬件（如 GPU）组成的云原生数据中心，以支撑日益庞大且多元的 AI 负载，包括离线训练任务、在线推理任务、
超参数优化（Hyperparameter Optimization, HPO）~\cite{rasley2017hyperdrive}任务以及强化学习（Reinforcement Learning, RL）智能体训练任务等~\cite{CAICT2024CloudWhitePaperCN}。
其中，离线训练与在线推理作为两类核心 AI 
计算负载，呈现出截然不同的系统特征与调度需求，各自催生了不同的调度范式~\cite{liqirui2022,luolei2021}。

针对离线训练任务，随着计算负载从传统深度学习模型转向大规模预训练模型，云数据中心的调度范式也随之发生了深刻变革。
早期以 ResNet~\cite{he2016deep}、BERT~\cite{devlin2019bert}等规模适中的模型为代表的分布式训练任务，
因其对计算资源具有“全有或全无（All-or-Nothing）”的分配需求，促使群组调度（Gang Scheduling）~\cite{gu2019tiresias,mahajan2020themis}成为该阶段的主流范式。
例如，Kubernetes~\cite{burns2016borg}等容器编排平台结合群组调度机制保障了训练作业的原子化启动。
此后，为打破这种刚性分配带来的资源利用率瓶颈，并适应新兴的弹性训练需求，集群调度技术逐渐向弹性调度范式~\cite{peng2018optimus,qiao2021pollux}发展。
Optimus~\cite{peng2018optimus}等代表性系统通过动态调整各训练作业的资源配额，在多租户环境下有效提升了
集群资源利用率与吞吐量。
近年来，随着 GPT-3~\cite{brown2020language}等超大参数量基础模型的兴起，训练任务对 GPU 算力、显存容量以及高速互联带宽的需求呈数量级攀升，
进一步对集群层面的高效资源调度提出了更加严苛的要求。

针对在线推理任务，在“模型即服务”（Model as a Service, MaaS）这一云端交付范式的持续驱动下，云数据中心的负载经历了从传统深度学习推理
向大语言模型（Large Language Model, LLM）推理，乃至 LLM 智能体（Agent）推理的快速演进。
在 MaaS 发展的早期，传统的深度学习推理请求通常相互独立且计算规模适中，常规的批处理调度即可满足服务需求。
然而，随着 LLM 推理服务逐渐成为在线推理的核心负载，其自回归生成的计算特性与庞大的访存开销
对端到端响应延迟与系统吞吐量提出了严格的服务质量（Quality of Service, QoS）挑战。
为此，现有的推理系统进行了大量底层架构优化，如 vLLM~\cite{kwon2023efficient}通过连续批处理（Continuous Batching）~\cite{yu2022orca}
与分页注意力（PagedAttention）机制大幅降低了显存碎片化开销，
SGLang~\cite{zheng2024sglang}则借助 RadixAttention 实现了跨请求的键值缓存（KV Cache）高效复用，
显著提升了单节点推理引擎的吞吐量。
然而，现代 LLM Agent 框架赋予了模型环境感知与多步规划能力~\cite{yao2023tree,yao2022react}，使得推理系统的服务对象由单次独立请求演变
为具有严格时序依赖的有向图（Directed Graph, DG）工作流。一次完整的 Agent 交互往往会触发大量子请求，这些子请求频繁复用共同的历史上下文，
从而对特定节点的 KV Cache 产生了强烈的局部性依赖~\cite{zheng2024sglang,lin2024parrot}，这为推理调度系统的设计带来了全新的机遇与挑战。

%除上述主流的离线训练与在线推理任务外，现代云数据中心还承载着超参数优化（Hyperparameter Optimization, HPO）~\cite{rasley2017hyperdrive}
%、强化学习（Reinforcement Learning, RL）智能体训练以及大规模数据预处理等多类特化任务，进一步加剧了集群负载的多样性与复杂性。
然而，无论是追求整体吞吐的大规模训练，还是具有严格延迟要求的在线推理，其高度动态的负载特征与多元化的性能优化目标，都给集群的资源调度带来了极高的系统复杂性。
传统面向大数据或高性能计算（HPC）的通用调度框架~\cite{wang2020survey}已难以有效适配此类高度复杂的应用场景。
因此，探索一种能够自适应 AI 负载特征、以全局性能为驱动的智能调度新范式，已成为打破算力瓶颈、
全面提升数据中心底层资源利用率与服务质量的核心所在。

%在多元化的 AI 负载生态中，深度学习离线分布式训练与基于大语言模型的多智能体在线推理，
%分别代表了当前大规模算力密集型任务与复杂逻辑依赖型业务的两类典型计算范式。前者作为基础模型开发与迭代的核心环节，
%针对该类任务的调度策略直接决定了集群算力资源的利用效率与模型收敛性能；
%后者作为当前大语言模型应用的主要形式，其执行性能直接关联到复杂交互场景下的服务质量与端到端响应时延。
%本文选取这两类在现代云数据中心最具代表性的核心负载作为主要研究对象，深入剖析其面临的调度挑战并提出针对性的系统优化方案，
%以此作为本研究的切入点。

\begin{figure}[htbp]
    \centering
    \includegraphics[width=0.8\textwidth]{figures/1.1.png}
    \caption{云数据中心中AI训推任务负载的调度架构示意图}
    \label{fig:datacenter}
\end{figure}

\section{研究意义}\label{sec:research-significance}
在多元化的 AI 负载生态中，深度学习离线分布式训练任务与基于大语言模型的多智能体在线推理任务，
分别代表了当前吞吐导向的离线批处理与延迟敏感的在线流式服务两类典型计算范式。前者作为基础模型开发与迭代的核心环节，
针对该类任务的调度策略直接决定了集群算力资源的利用效率与模型收敛性能；
后者作为当前大语言模型应用的主要形式，其执行性能直接关联到复杂交互场景下的服务质量与端到端响应时延。
本文选取这两类在现代云数据中心最具代表性的 AI 负载作为主要研究对象，深入剖析其面临的调度挑战，
并提出针对性的调度优化方案。

在深度学习训练任务调度领域，尽管当前的弹性调度范式在提升集群资源利用率上取得了一定进展，但其底层决策机制的局限性正日益凸显：
随着大规模模型训练需求的演进，现代集群调度需要综合考量的系统约束与资源配置不断增多，
例如底层硬件的异构性、多维并行策略的动态配置，以及复杂的任务依赖关系等。
上述多种因素的深度耦合，导致集群调度的决策空间发生指数级膨胀。
面对庞大且复杂的解空间，现有调度系统所依赖的启发式规则或数学规划求解器陷入了固有瓶颈：
启发式算法在面对高维、非线性的异构集群状态时建模能力不足，难以求得最优的分配策略；
而精确求解器则极易陷入维度灾难，无法满足在线动态调度的实时性要求~\cite{qiao2021pollux,chen2023deepboot}。
因此，迫切需要引入具备高维状态感知与泛化能力的自适应决策机制，在保证在线调度实时性的同时逼近全局最优解，
以实现高效的集群资源分配。

%与离线训练侧重利用率与收敛性不同，随着生成式大模型的迅速发展，以"模型即服务"（Model as a Service, MaaS）为代表的
%在线推理服务逐渐成为云数据中心的另一类核心计算负载，其调度核心也转向了端到端响应延迟与系统吞吐量~\cite{luowen2024}。
%为了应对 LLM 推理对低延迟与高吞吐的严格要求，推理引擎在系统层面经历了持续的技术迭代~\cite{gexuran2025,zhou2024surveyefficientinferencelarge,yi2025edgemoeempoweringsparselarge}。
%vLLM~\cite{kwon2023efficient}、SGLang~\cite{zheng2024sglang} 等高性能推理框架
%通过引入连续批处理（Continuous Batching）~\cite{yu2022orca}与分页注意力（PagedAttention）~\cite{kwon2023efficient}
%等机制，有效缓解了推理过程中的显存碎片化问题，显著提升了内存利用率。

另一方面，针对在线推理服务，尤其是前文所述的 LLM Agent 复杂有向图工作流，现有的系统架构暴露出了严重的调度缺陷。
尽管单节点推理引擎（如 vLLM~\cite{kwon2023efficient} 等）在底层显存管理上取得了显著进展，
但仅凭实例级的优化已难以应对集群级别的复杂调度挑战。
在分布式架构下，现有的轮询或会话亲和性（Session Affinity）等全局路由策略，本质上仍建立在“请求相互独立”的假设之上。
它们完全无法感知 Agent 子请求之间复杂的拓扑依赖，也无法获取全局分散的键值缓存分布状态，
极易导致具有强上下文依赖的并发任务被错误地打散至不同节点，进而引发严重的队头阻塞
与端到端响应的长尾延迟~\cite{zheng2024sglang,lin2024parrot}。
因此，必须专门设计具备全局状态感知与复杂工作流并发处理能力的调度机制。

%然而，单节点引擎层面的优化难以应对LLM Agent 应用带来的系统级挑战。
%与早期 LLM 服务中单向、连续的文本生成不同，现代 Agent 框架要求 LLM 具备环境感知、
%工具调用及多步规划能力~\cite{yao2023tree,yao2022react}。
%因此，LLM推理系统所服务的对象已经由单次独立的调用请求，转变为包含中间状态生成、环境反馈与动态分支的复杂工作流。

%在系统负载抽象层面，Agent 请求呈现出具有严格时序依赖的有向图（Directed Graph, DG）工作流结构，
%打破了传统LLM服务中“请求相互独立”的固有假设。一次完整的 Agent 交互通常包含大量细粒度的并发推理子请求，
%由于这些子请求频繁复用共同的历史上下文，此类工作流在执行过程中会对特定推理实例上的 KV Cache 状态产生局部性依赖~\cite{zheng2024sglang,lin2024parrot}。

%上述复杂工作流形态的大量涌现，使得基于独立请求假设的传统集群调度抽象面临严峻挑战。在分布式推理架构中，
%随机散列或轮询等负载均衡策略无法感知子请求间的拓扑依赖，
%也无法联合建模全局分散的键值缓存（KV Cache）状态与节点队列负载。
%这种缺乏全局感知的调度机制会导致具有强依赖关系的推理任务被盲目打散至不同的节点，
%引发严重的队头阻塞，进而导致端到端请求的长尾延迟显著增加。因此，为应对 LLM Agent 在复杂并
%发场景下的调度需求，需专门设计具备全局感知能力的请求路由策略与并发调度框架。
综上所述，无论是旨在提升集群吞吐量的离线深度学习训练，还是追求低延迟的在线 LLM Agent 推理，
现有的通用调度范式均难以满足其严苛的服务需求。
立足于这一现实背景，本文针对这两类最具代表性的核心 AI 负载，分别设计并提出了定制化的集群调度算法，
旨在为提升云原生 AI 算力基础设施的资源利用率与服务质量提供可行的解决方案。

\section{主要工作与创新点}

针对前文所述的系统调度痛点，本文的核心思想是引入深度强化学习（Deep Reinforcement Learning, DRL）作为统一的求解与决策框架。
正如第\ref{sec:research-significance}节所述，现代 GPU 集群调度普遍面临着难以兼顾
求解质量与响应速度的固有矛盾。
而深度强化学习范式为打破这一瓶颈提供了可行途径：
其基于深度神经网络的强表征能力，能够联合拟合多维底层资源指
标与复杂的任务依赖关系，突破启发式规则的建模局限以逼近全局最优；更重要的是，在部署推理时，
调度器仅需执行多次神经网络的前向传播即可得到完整的调度策略，
从而在保障决策质量的同时实现了极致的响应速度。

基于上述统一的 DRL 求解思路，本文分别针对 DL 离线训练任务与 LLM Agent 在线推理任务两大核心负载，
设计并实现了两套基于强化学习的调度算法。
考虑到这两类调度问题在样本特性与决策结构上存在本质差异，本文分别采用 on-policy 的
REINFORCE 算法与 off-policy 的 Double DQN 算法进行针对性求解。
具体而言，DL 训练任务通常持续时间较长，且伴随着其运行期间集群状态的高度动态变化，
系统能够在连续的调度周期中积累充足的时序交互样本；
此外，训练集群的资源分配本质上是一个多步序列化的决策过程，中间步骤的分配动作与最终的全局奖励信号之间存在强耦合。
基于上述两点特性，本文采用 on-policy 的 REINFORCE 算法，以确保策略梯度估计的无偏性。
相比之下，在 LLM Agent 推理路由场景中，每条请求的调度决策相对独立且仅触发单次动作，
可用的训练样本规模严格受限于在线请求的到达量。为克服采样效率较低的问题，
本文选用 off-policy 的 Double DQN 算法，旨在借助经验回放
机制充分复用历史决策样本，从而显著提升样本利用效率。

基于上述考量，本文的第\ref{chap:adasched}章聚焦深度学习离线训练场景，提出并实现了性能驱动的弹性集群调度算法 AdaSched；
第\ref{chap:agent-routing}章聚焦 LLM Agent 在线推理场景，提出了融合动态感知排队与强化学习分发的全局路由算法 RAGS。

在第\ref{chap:adasched}章，本文聚焦深度学习离线训练场景，
提出以模型验证性能为截止判据的弹性集群调度算法 AdaSched。
该算法将调度回报定义为集群内作业性能目标满足程度的加权累积得分，
支持用户为每个训练任务提交多级性能目标（每级对应不同资源预算），并将集群调度形式化为最大化累积回报的
整数线性规划问题。鉴于该问题具有极高的状态空间复杂度，传统精确求解器难以满足动态集群在线调度的实时性要求，
本文引入深度强化学习方法来高效逼近最优调度策略。在实现上，
AdaSched 采用 Actor-Critic~\cite{sutton1999policy} 架构，其策略网络创新性地融合
了门控循环单元（GRU）~\cite{cho2014learning}与双级注意力机制，深度联合感知多步调度的时序依赖
与全局集群上下文；同时，设计了阶段式密集奖励函数以有效缓解稀疏奖励问题，并辅以动作掩码机制严格
确保调度决策的合法性。基于 Pollux 数据集与 Philly~\cite{jeon2019analysis} 负载轨迹的仿真实验表明，
AdaSched 的平均调度回报分别达到弹性调度器 ElasticFlow 的 1.2 倍与非弹性调度器 Chronus 的 2.76 倍，
且调度决策延迟降至 ElasticFlow 的约五分之一，在训练集群场景下兼顾了调度质量与响应效率。

在第\ref{chap:agent-routing}章，本文针对 LLM Agent 在线推理场景，
提出了面向 LLM Agent 服务系统的全局路由算法 RAGS（RL-based Agent Global Scheduler）。
该算法包含动态感知排队算法 DynA 与基于强化学习的全局分发策略两个核心组件。具体而言，DynA 无需
预先获取 Agent 的工作流结构，而是通过在线统计各 Agent 类型工作流的完成轮次概率分布，以条件近完成概率
作为跨 Agent 调度优先级；同时，结合累积延迟驱动的 Agent 内排序机制与基于排队年龄的防饥饿策略，
有效兼顾了系统的延迟与公平性。强化学习分发策略则将分发决策建模为马尔可夫决策过程。
其状态空间联合表征了两类核心信息：一是待调度请求的输入输出等特征，二是各 LLM 实例的队列负载、显存占
用率及 KV Cache 命中率等底层运行状态。
该策略采用基于 Transformer 的网络架构，引入类型嵌入区分请求与实例特征；同时设计了融合 KV 缓存复用、负载均衡与系统吞吐量的复合奖励函数，
通过离线预训练与在线微调的两阶段流程持续适配动态负载。
实验结果表明，在多轮对话、文档摘要与代码生成均匀混合的工作负载下，RAGS 相较 Ayo 与 Parrot 分别将平均端到端延迟降低了 $14.5\%$ 与 $28.6\%$；
消融实验进一步验证了基于强化学习的分发策略与 DynA 排队策略在高负载下分别贡献
约 42.27\% 与 5.97\% 的延迟优化，证明了RAGS算法的有效性。

综上所述，本文的主要创新点如下：
\begin{itemize}
    \item 提出并实现面向深度学习训练集群的性能驱动弹性调度算法 AdaSched。该方法创新性地引入模型性能替代用户
          预估迭代次数作为调度截止判据，并设计了多级性能目标机制将集群调度形式化为最大化累积调度回报问题；构建了融合
          GRU 时序编码与两级注意力机制的上下文感知策略网络，精准评估各任务的调度紧迫程度；通过动作掩码
          与阶段式密集奖励设计提升强化学习训练稳定性与决策合法性。
    \item 提出并实现面向 LLM Agent 服务系统的全局路由算法 RAGS，包含 DynA 排队算法与基于强化学习的分发策略。
          DynA 基于在线完成轮次分布统计实现无需预知工作流结构的工作流进度感知优先排队；
          分发策略基于 Transformer 策略网络统一表征多维请求与实例特征实现全局感知的分发决策。
          设计了一种融合 KV 缓存复用、负载均衡与吞吐量代理的复合奖励函数，并采用离线预训练与在线微调
          的两阶段训练范式，使路由策略能够适应动态负载。
\end{itemize}

\section{本文结构与章节安排}

本文共分为五章，各章节的内容安排如下：

第一章为绪论。该章节首先阐述本文的研究背景。
其次，深入剖析云数据中心在面对深度学习离线训练与 LLM Agent 在线推理这两类核心负载时所面临的调度挑战，进而引出本文的研究意义。
在此基础上，提炼并总结本文的主要工作与创新点，最后阐述全文的结构安排。  

%第二章为相关技术与研究现状。该章节主要介绍本文所涉及的理论知识与相关领域的研究现状，涵盖深度学习集群调度框架、大语言模型推理与路由技术以及深度强化学习算法。
%在此基础上，梳理国内外相关领域的最新研究进展，指出现有调度系统与路由策略在应对复杂系统状态与严格时序依赖时的局限性，为本文后续提出的方法奠定基础。
第二章为背景知识与研究现状。该章节系统介绍本文所涉及的理论知识，并在此基础上梳理国内外在 DL 训练集群与LLM 推理集群调度的最新研究进展，
指出现有调度系统与路由策略在应对复杂系统状态与严格时序依赖时的局限性，为本文工作奠定基础。

第三章为基于深度强化学习的性能驱动集群调度方法。该章节首先分析当前弹性调度方法以迭代次数作为判据导致的资源浪费与准入控制失效问题。
继而提出以模型验证性能为截止判据的 AdaSched 调度算法，并详细阐述策略网络设计，包括门控循环单元与两级注意力机制的应用。
最后通过仿真实验对比分析，验证 AdaSched 在提升系统调度回报与响应效率上的有效性。

第四章为基于强化学习的 LLM Agent 服务系统全局调度算法。
本章先分析 Agent 工作流中请求的拓扑依赖与异构性特征导致的队头阻塞与长尾延迟问题。
然后详细介绍所提出的全局调度算法 RAGS 中动态感知排队算法 DynA 与基于强化学习的全局路由算法这两个
核心组件的设计与实现机制。最后经过一系列的多轮对话、代码生成等混合工作负载的对比与消融实验，
验证了 RAGS 在降低端到端长尾延迟方面的可行性与优越性。

第五章为总结与展望。该章节对本文的两个主要研究工作进行全面总结，反思并讨论了当前方法在更通用场景或更大规模集群下存在的局限性与不足，
最后展望了本文工作可能的改进方向。

\cleardoublepage
